\chapter{Переприоритизация программы после заморозки компакттона}
\label{ch:v6-spectral-transition-post-compacton-program-reprioritization-gate}

\section{Почему требуется смена роли, а не новая поправка}

Статусная заморозка доказала две разные вещи. Во-первых, дискретная
составная модель содержит точное локализованное решение и вычислимый
радиационный форм-фактор. Во-вторых, это решение не является автономно
рождающейся частицей: нет общего trigger, физической скорости,
единственного устойчивого endpoint и абсолютного масштаба.

Следовательно, запрещены два симметричных обхода. Нельзя вновь объявить
compacton частицей после добавления ещё одного граничного коэффициента.
Но нельзя и выбросить его точные данные только потому, что первоначальная
физическая роль оказалась неверной. Требуется проверить, имеют ли они
другую функцию внутри уже построенной проекторной динамики.

Настоящий гейт не выводит новый механизм. Он сравнивает допустимые маршруты
и выбирает один следующий falsification test.

\section{Два независимых незамкнутых ledger}

\subsection{Проекторная кристаллизация}

Для семейного параметра порядка строго известны:
\begin{equation}
 \beta_c=1.5426695409\ldots,
 \qquad
 \beta_{\rm sp}=\frac{21}{2},
 \label{eq:v6-post-compacton-projective-thresholds}
\end{equation}
а также энтропийный и энергетический бюджеты перехода к состоянию
сосуществования,
\begin{equation}
 \Delta S=0.7402345698\ldots,
 \qquad
 \Delta E=0.4798400112\ldots.
 \label{eq:v6-post-compacton-ordering-budget}
\end{equation}
Четырёхсостояний носитель имеет достаточную энтропийную ёмкость
$\log4$, но невырожденные четырёхтактные часы не являются автономным
холодильником: их точный энергетический потолок выделенного веса равен
$2/3$, тогда как требуется $0.9121665963\ldots$.

Аффинный угол $P_3$ содержит настоящую резонансную кратность три, однако
каноническая связь $\rho V$ изотропна и не создаёт расщепление $1+2$.
Таким образом, статическая поверхность перехода и вместимость стока
известны, но закон внутреннего охлаждения $\beta(\tau)$ не выведен.

\subsection{Дискретная радиация}

Для эта-ориентированной характерной квадратуры получены
\begin{equation}
 \rho_{\rm rad}(0)=2\pi,
 \qquad
 \Gamma_{\rm cycle}=4\pi^2|\delta|^2.
 \label{eq:v6-post-compacton-radiation-law}
\end{equation}
Этот канал не стабилизирует compacton: его ядро теряет норму, а
нежелательный характер не затухает. Но сам факт исходящего переноса в
пространственно разделяющиеся моды остаётся точным. В отличие от конечных
четырёхтактных часов, последовательность исходящих пакетов не рекуррентна
в проверенном до оборачивания объёме и естественно хранит запись о
прошедших циклах.

Открытый вопрос поэтому меняется. Требуется не ``может ли излучение
выбрать compacton'', а ``может ли тот же исходящий канал быть частью
замкнутого энергетического баланса проекторной кристаллизации''.

\section{Реестр маршрутов}

Машинный аудит сравнивает шесть вариантов по семи критериям: соблюдение
заморозки, отказ от выбора одной ветви $\pm i$, повторное использование
выведенных данных, атака на дефицит охлаждения, атака на R2 или R3,
отсутствие нового размерного входа на стадии выбора и наличие резкого
следующего теста.

\begin{center}
\begin{tabular}{p{0.43\textwidth}p{0.14\textwidth}p{0.31\textwidth}}
\toprule
Маршрут & Счёт & Решение \\
\midrule
Вернуть compacton как endpoint & $2/7$ & запрещён freeze-гейтом \\
Вставить jump или выбрать одну ветвь & $1/7$ & ручной $\gamma$ и ось \\
Сначала назначить абсолютный масштаб & $3/7$ & динамически преждевременно \\
Начать произвольный новый родитель & $5/7$ & допустимо, но не минимально \\
Real-парный радиационный мост охлаждения & $7/7$ & следующий аудит \\
Завершить Том VI без последнего моста & $3/7$ & резерв при провале \\
\bottomrule
\end{tabular}
\end{center}

Выбор пятого маршрута не является положительным результатом. Он означает
только, что этот вопрос одновременно использует два строгих ledger и
может быть закрыт одним вычислимым стоп-критерием.

\section{Новая карта ролей}

Следующий тест обязан использовать четыре объекта с разными функциями:
\begin{enumerate}
 \item проекторное состояние $R$ --- охлаждаемая система и параметр
 порядка;
 \item глобально нейтральная Real-пара --- допустимая структура перехода,
 без требования выбрать одну сопряжённую ветвь;
 \item compacton --- вспомогательный конечномерный осциллятор и источник
 уже измеренного форм-фактора, но не конечная частица;
 \item исходящие пространственные моды --- кандидат переноса энергии и
 энтропии, но не заранее объявленный тепловой резервуар.
\end{enumerate}

Такое назначение согласуется с общим принципом автономных квантовых машин:
часы и сток являются физическими ресурсами и испытывают обратную реакцию,
а не задают внешний параметр времени \cite{post-compacton-erker2017}.
Минимальные автономные холодильники также требуют резонансной структуры,
а не одной энтропийной вместимости
\cite{post-compacton-linden2010}. Исходящий непрерывный канал может
создавать необратимость только после проверки знака потока и общего
гамильтониана \cite{post-compacton-fano1961}.

\section{Контракт следующего родителя}

Гейт
\path{version6_spectral_transition_real_pair_radiative_cooling_parent_gate}
должен вывести, а не предположить, шесть утверждений.
\begin{enumerate}
 \item Суммарный по Real-паре радиационный ток должен быть чётным и не
 сокращаться в полном KO6-завершении.
 \item Должен существовать один закон сохранения энергии для $R$,
 compacton-осциллятора и исходящих мод.
 \item Амплитуда $\delta$ в
 \eqref{eq:v6-post-compacton-radiation-law} должна следовать из состояния
 или родительской связи, а не быть новым начальным параметром.
 \item Из замкнутой динамики должен следовать монотонный редуцированный
 параметр, эквивалентный $\beta(\tau)$, без внешнего quench.
 \item Для ненулевого множества допустимых начальных состояний система
 должна пересечь $\beta_c$ или $\beta_{\rm sp}$.
 \item Вспомогательный источник должен существовать достаточно долго,
 чтобы завершить перенос, а не разрушиться раньше порога.
\end{enumerate}

\begin{nogocriterion}[Запрет переименования излучения в охлаждение]
Ненулевая норма исходящего пакета ещё не является тепловым потоком.
Маршрут немедленно закрывается, если Real-половины сокращают ток, перевод
радиационной нормы в энергию требует нового коэффициента, $\delta$ или
$\beta(\tau)$ задаются вручную, либо общий энергетический баланс не
существует.
\end{nogocriterion}

Особенно важно не использовать выбранную $C_4$-ось как скрытое начальное
условие. Родитель должен быть сформулирован на Real-чётной паре; если
ненулевой ток появляется только после выбора одной характерной линии,
тест считается проваленным.

\begin{theorem}[Переприоритизация после compacton]
Заморозка compacton как частицы не запрещает использовать его доказанный
радиационный форм-фактор как вход нового родительского теста. Среди шести
рассмотренных маршрутов единственным минимальным вариантом, который
соблюдает freeze-гейт и одновременно атакует незамкнутый закон внутреннего
охлаждения, является Real-парный радиационный мост к проекторному переходу.
Это выбор следующего falsification test, а не доказательство механизма.
\end{theorem}

\section{Нулевая гипотеза и стоп-критерий}

Нулевая гипотеза: полный Real/KO6 ток либо сокращается, либо его связь с
проектной энергией содержит новый свободный коэффициент; тогда точная
радиация не задаёт внутреннего охлаждения.

Стоп-критерий: достаточно одного из пяти исходов --- сокращение парного
тока, свободная конверсионная нормировка, внешний $\delta$ или
$\beta(\tau)$, отсутствие общего сохранения энергии либо разрушение
источника до пересечения $\beta_c$. При выполнении стоп-критерия Том VI
завершается без нового динамического механизма.

\begin{protocol}
\begin{enumerate}
 \item сравнённых маршрутов: \textbf{шесть};
 \item критериев выбора: \textbf{семь};
 \item compacton как частица: \textbf{не переоткрыт};
 \item выбор одной ветви $\pm i$: \textbf{не требуется};
 \item проекторная спинодаль: \textbf{$21/2$};
 \item необходимый экспорт энтропии: \textbf{$0.7402345698\ldots$};
 \item радиационная плотность: \textbf{$2\pi$};
 \item коэффициент за цикл: \textbf{$4\pi^2|\delta|^2$};
 \item закон внутреннего охлаждения: \textbf{ещё не выведен};
 \item выбранный маршрут: \textbf{Real-парный радиационный мост};
 \item статус маршрута: \textbf{только следующий строгий аудит};
 \item резерв при провале: \textbf{завершение Тома VI без модели рождения}.
\end{enumerate}
\end{protocol}

Машинный сертификат сохранён в
\path{s2t_v6_spectral_transition_post_compacton_program_reprioritization_gate_results.json}.

\begin{thebibliography}{9}
\bibitem{post-compacton-erker2017}
P.~Erker, M.~T.~Mitchison, R.~Silva, M.~P.~Woods, N.~Brunner, and
M.~Huber, \emph{Autonomous Quantum Clocks: Does Thermodynamics Limit Our
Ability to Measure Time?}, Physical Review X 7 (2017), 031022,
\path{arXiv:1609.06704}.

\bibitem{post-compacton-linden2010}
N.~Linden, S.~Popescu, and P.~Skrzypczyk,
\emph{How Small Can Thermal Machines Be? The Smallest Possible
Refrigerator}, Physical Review Letters 105 (2010), 130401,
\path{doi:10.1103/PhysRevLett.105.130401}.

\bibitem{post-compacton-fano1961}
U.~Fano,
\emph{Effects of Configuration Interaction on Intensities and Phase Shifts},
Physical Review 124 (1961), 1866--1878,
\path{doi:10.1103/PhysRev.124.1866}.
\end{thebibliography}